\cotefile{36}
\scanpage{53}
\mspage{-1.6-}

$\phi:V_1\longrightarrow V_2$ commuting with the dual actions, then $\phi\circ E_1=E_2\circ\phi$ where $E_1$ and $E_2$ are the Reynolds operators on $V_1$ and $V_2$ respectively.

Suppose $V$ is actually a ring $R$, and $\hat\sigma$ is an action on the ring. Then the image of $E$ is the subspace of invariant elements of $R$, which is a \uline{subring} $R_0$ of $R$. Although $E$ is not a ring homomorphism, it satisfies the well-known Reynolds identity:
\[
(*)\qquad\text{If }x\in R,\ y\in R_0,\text{ then }E(x.y)=(Ex)\cdot y.
\]
(cf. Cartier [6], 7-08). To prove this, set $R_1=\ker(E)$. Then $R\cong R_0\oplus R_1$, and $(*)$ is equivalent to the 2 statements $R_0\cdot R_0\subset R_0$ and $R_0\cdot R_1\subset R_1$. To prove the latter, let $g\in R_0$, and let $V\subset R_1$ be a (finite dimensional) irreducible invariant subspace. Then either $g\cdot V=(0)$ or multiplication by $g$ sets up an isomorphism of $V$ with an irreducible invariant subspace $g\cdot V$. Since $g$ is an invariant, the representations of $G$ on $V$ and $g\cdot V$ are isomorphic and non-trivial. Therefore $g\cdot V\subset R_1$.

\fbox{\S2.} With these preliminaries disposed of, we can now take up the first case of the theory - the action of a reductive algebraic group on an affine scheme $X$. To keep straight the techniques in the proof, we do not assume that $X$ is an algebraic scheme over $k$.

\uline{Theorem 1.1}: Let $X$ be an affine scheme over $k$, let $G$ be a reductive algebraic group, and let $\sigma:G\mathbin{\times}X\longrightarrow X$ be an action of $G$ on $X$. Then a universal categorical quotient $(Y,\phi)$ of $X$ by $G$

\margin{\unsure{52}}
